% GENERATED FILE. Edit canonical lessons, then run npm run generate:books.
% canonical-source-hash: fnv1a64:42b6259e8112f80a
% canonical-lessons: LA-C380-conduco, LA-C380-loco, LA-C380-famulus, LA-C380-navo, LA-C380-fatigo, LA-R380-review-373-380

\chapter{Hiring and Hard Labour}
\label{ch:la-hiring-and-hard-labour}
\hlchaptermodality{380}

\begin{chapteropening}
\textbf{By the end of this chapter:} \emph{I can say that someone hires workers or lets out a house, speak of a servant, and say that someone works hard or wears himself out.}

Complete the last lesson of chapter 380: I can say that someone hires workers or lets out a house, speak of a servant, and say that someone works hard or wears himself out.
\end{chapteropening}

\section[condūcō, condūcere]{condūcō, condūcere — to hire, to rent; to bring together}

\label{lesson:LA-C380-conduco}



\begin{quote}
Before the new one: say the Latin for toilsome, then the Latin for care.
\end{quote}

\subsection*{You'll want to know: condūcō, condūcere}
\textbf{condūcō, condūcere} — \textquotedblleft{}to hire, to rent; to bring together\textquotedblright{}.

\begin{quote}
Agricola operāriōs condūxit.
\end{quote}

\begin{quote}
\emph{The farmer hired labourers.}
\end{quote}

\emph{condūcō, condūcere, condūxī, conductum}. The man who hires or rents is the \emph{conductor}. English \emph{conduct}.

\begin{cousinweb}
\emph{condūcō} is built from \emph{con-} \textquotedblleft{}together\textquotedblright{} and \emph{dūcō} \textquotedblleft{}to lead\textquotedblright{} (EtymDB; Lewis \& Short): to lead workers together. English \emph{conduct} comes from it.
\end{cousinweb}

\begin{grammarlens}[title={What you've built}]
One more word for telling, explaining and describing at B1.
\end{grammarlens}

\subsection*{Your turn}
\begin{itemize}
\raggedright
  \item \emph{Say it:} \emph{condūcō, condūcere}
  \item \emph{Say it:} \emph{condūcō, condūcere}, once more
  \item \emph{Say it:} a sentence of your own with \emph{condūcō, condūcere}
  \item \emph{Recall:} say the Latin for toilsome, then the Latin for care, then say \emph{condūcō, condūcere} again
\end{itemize}

\begin{tcolorbox}[breakable,colback=teal!4,colframe=teal!35!black,title={Before you move on}]
What does condūcō, condūcere mean? (\textquotedblleft{}To hire, to rent; to bring together\textquotedblright{}.) Say it once more.
\end{tcolorbox}
\section[locō, locāre]{locō, locāre — to place; to let out, to hire out}

\label{lesson:LA-C380-loco}



\begin{quote}
Before the new one: say the Latin for care, then the Latin for to hire.
\end{quote}

\subsection*{You'll want to know: locō, locāre}
\textbf{locō, locāre} — \textquotedblleft{}to place; to let out, to hire out\textquotedblright{}.

\begin{quote}
Dominus domum suam mercēde magnā locāvit.
\end{quote}

\begin{quote}
\emph{The owner let his house for a high rent.}
\end{quote}

\emph{locō, locāre, locāvī, locātum}, built on \emph{locus}. In Roman law \emph{locāre} is to let out and \emph{condūcere} to take on hire. English \emph{locate}.

\begin{cousinweb}
\emph{locō} is built on \emph{locus} \textquotedblleft{}a place\textquotedblright{} (EtymDB; Lewis \& Short): to put in place, then to let out for hire. English \emph{locate} comes from it.
\end{cousinweb}

\begin{grammarlens}[title={What you've built}]
One more word for telling, explaining and describing at B1.
\end{grammarlens}

\subsection*{Your turn}
\begin{itemize}
\raggedright
  \item \emph{Say it:} \emph{locō, locāre}
  \item \emph{Say it:} \emph{locō, locāre}, once more
  \item \emph{Say it:} a sentence of your own with \emph{locō, locāre}
  \item \emph{Recall:} say the Latin for care, then the Latin for to hire, then say \emph{locō, locāre} again
\end{itemize}

\begin{tcolorbox}[breakable,colback=teal!4,colframe=teal!35!black,title={Before you move on}]
What does locō, locāre mean? (\textquotedblleft{}To place; to let out, to hire out\textquotedblright{}.) Say it once more.
\end{tcolorbox}
\section[famulus, famulī]{famulus, famulī — a household servant, an attendant}

\label{lesson:LA-C380-famulus}



\begin{quote}
Before the new one: say the Latin for to hire, then the Latin for to place.
\end{quote}

\subsection*{You'll want to know: famulus, famulī}
\textbf{famulus, famulī} — \textquotedblleft{}a household servant, an attendant\textquotedblright{}.

\begin{quote}
Famulus cēnam in mēnsā posuit.
\end{quote}

\begin{quote}
\emph{The servant laid dinner on a table.}
\end{quote}

\emph{famulus, famulī} m. \emph{Familia} first meant the household's servants, the \emph{famulī}. English \emph{family} and \emph{familiar} come from that word.

\begin{cousinweb}
\emph{famulus} is inherited from Proto-Italic \emph{*famelos} (EtymDB), and \emph{familia} \textquotedblleft{}a household\textquotedblright{} is built on it (Lewis \& Short). English \emph{family} and \emph{familiar} come from \emph{familia}.
\end{cousinweb}

\begin{grammarlens}[title={What you've built}]
One more word for telling, explaining and describing at B1.
\end{grammarlens}

\subsection*{Your turn}
\begin{itemize}
\raggedright
  \item \emph{Say it:} \emph{famulus, famulī}
  \item \emph{Say it:} \emph{famulus, famulī}, once more
  \item \emph{Say it:} a sentence of your own with \emph{famulus, famulī}
  \item \emph{Recall:} say the Latin for to hire, then the Latin for to place, then say \emph{famulus, famulī} again
\end{itemize}

\begin{tcolorbox}[breakable,colback=teal!4,colframe=teal!35!black,title={Before you move on}]
What does famulus, famulī mean? (\textquotedblleft{}A household servant, an attendant\textquotedblright{}.) Say it once more.
\end{tcolorbox}
\section[nāvō, nāvāre]{nāvō, nāvāre — to do with energy (operam nāvāre: to work hard)}

\label{lesson:LA-C380-navo}



\begin{quote}
Before the new one: say the Latin for to place, then the Latin for a household servant.
\end{quote}

\subsection*{You'll want to know: nāvō, nāvāre}
\textbf{nāvō, nāvāre} — \textquotedblleft{}to do with energy (operam nāvāre: to work hard)\textquotedblright{}.

\begin{quote}
Amīcō operam nāvāvit.
\end{quote}

\begin{quote}
\emph{He worked hard for his friend.}
\end{quote}

\emph{nāvō, nāvāre, nāvāvī, nāvātum}, built on \emph{nāvus} (old \emph{gnāvus}), busy; \emph{ignāvus}, lazy, is the opposite. It mostly appears in \emph{operam nāvāre} + dative.

\begin{cousinweb}
\emph{nāvō} is built on \emph{nāvus}, old \emph{gnāvus} \textquotedblleft{}busy, active\textquotedblright{} (EtymDB; Lewis \& Short), the adjective that \emph{ignāvus} \textquotedblleft{}lazy\textquotedblright{} denies.
\end{cousinweb}

\begin{grammarlens}[title={What you've built}]
One more word for telling, explaining and describing at B1.
\end{grammarlens}

\subsection*{Your turn}
\begin{itemize}
\raggedright
  \item \emph{Say it:} \emph{nāvō, nāvāre}
  \item \emph{Say it:} \emph{nāvō, nāvāre}, once more
  \item \emph{Say it:} a sentence of your own with \emph{nāvō, nāvāre}
  \item \emph{Recall:} say the Latin for to place, then the Latin for a household servant, then say \emph{nāvō, nāvāre} again
\end{itemize}

\begin{tcolorbox}[breakable,colback=teal!4,colframe=teal!35!black,title={Before you move on}]
What does nāvō, nāvāre mean? (\textquotedblleft{}To do with energy (operam nāvāre: to work hard)\textquotedblright{}.) Say it once more.
\end{tcolorbox}
\section[fatīgō, fatīgāre]{fatīgō, fatīgāre — to tire out, to weary}

\label{lesson:LA-C380-fatigo}



\begin{quote}
Before the new one: say the Latin for a household servant, then the Latin for to do with energy.
\end{quote}

\subsection*{You'll want to know: fatīgō, fatīgāre}
\textbf{fatīgō, fatīgāre} — \textquotedblleft{}to tire out, to weary\textquotedblright{}.

\begin{quote}
Longum iter equōs fatīgāvit.
\end{quote}

\begin{quote}
\emph{The long journey tired out the horses.}
\end{quote}

\emph{fatīgō, fatīgāre, fatīgāvī, fatīgātum}. English \emph{fatigue}.

\begin{cousinweb}
\emph{fatīgō} is linked by Lewis \& Short with \emph{affatim} \textquotedblleft{}enough\textquotedblright{} and \emph{fatīscō} \textquotedblleft{}to grow weak\textquotedblright{}, but its full origin is not settled. English \emph{fatigue} comes from it through French.
\end{cousinweb}

\begin{grammarlens}[title={What you've built}]
One more word for telling, explaining and describing at B1.
\end{grammarlens}

\subsection*{Your turn}
\begin{itemize}
\raggedright
  \item \emph{Say it:} \emph{fatīgō, fatīgāre}
  \item \emph{Say it:} \emph{fatīgō, fatīgāre}, once more
  \item \emph{Say it:} a sentence of your own with \emph{fatīgō, fatīgāre}
  \item \emph{Recall:} say the Latin for a household servant, then the Latin for to do with energy, then say \emph{fatīgō, fatīgāre} again
\end{itemize}

\begin{tcolorbox}[breakable,colback=teal!4,colframe=teal!35!black,title={Before you move on}]
What does fatīgō, fatīgāre mean? (\textquotedblleft{}To tire out, to weary\textquotedblright{}.) Say it once more.
\end{tcolorbox}
\section[(dialogue)]{Words from chapters 373–380 — a review}

\label{lesson:LA-R380-review-373-380}



\begin{quote}
Say the Latin for a result and to follow.
\end{quote}

\subsection*{Your turn}
Say these aloud:

\begin{itemize}
\raggedright
  \item a result, to follow, to come forth, to turn out well, to overflow
  \item to harm, to be useful, to be a hindrance, to hinder, to throw into confusion
  \item to touch, to arise, by chance, unexpected, a movement
  \item a work, a craftsman, an artist, a labourer, to work
  \item a fuller, a goldsmith, a carpenter, a mason, a sculptor of statues
  \item to train, to make, to weave, to sew, to carve
  \item gain, pay, a hired worker, toilsome, care
  \item to hire, to place, a household servant, to do with energy, to tire out
\end{itemize}

\begin{tcolorbox}[breakable,colback=teal!4,colframe=teal!35!black,title={Before you move on}]
A result? (\textbf{effectus, effectūs}.) To follow? (\textbf{cōnsequor, cōnsequī}.) To come forth? (\textbf{prōveniō, prōvenīre}.) To turn out well? (\textbf{succēdō, succēdere}.) To overflow? (\textbf{redundō, redundāre}.) To harm? (\textbf{noceō, nocēre}.) To be useful? (\textbf{prōsum, prōdesse}.) To be a hindrance? (\textbf{obsum, obesse}.) To hinder? (\textbf{impediō, impedīre}.) To throw into confusion? (\textbf{perturbō, perturbāre}.) To touch? (\textbf{contingō, contingere}.) To arise? (\textbf{exsistō, exsistere}.) By chance? (\textbf{fortuītus, fortuīta, fortuītum}.) Unexpected? (\textbf{inopīnātus, inopīnāta, inopīnātum}.) A movement? (\textbf{mōtus, mōtūs}.) A work? (\textbf{opus, operis}.) A craftsman? (\textbf{opifex, opificis}.) An artist? (\textbf{artifex, artificis}.) A labourer? (\textbf{operārius, operāriī}.) To work? (\textbf{operor, operārī}.) A fuller? (\textbf{fullō, fullōnis}.) A goldsmith? (\textbf{aurifex, aurificis}.) A carpenter? (\textbf{lignārius, lignāriī}.) A mason? (\textbf{caementārius, caementāriī}.) A sculptor of statues? (\textbf{statuārius, statuāriī}.) To train? (\textbf{exerceō, exercēre}.) To make? (\textbf{fabricor, fabricārī}.) To weave? (\textbf{texō, texere}.) To sew? (\textbf{suō, suere}.) To carve? (\textbf{sculpō, sculpere}.) Gain? (\textbf{quaestus, quaestūs}.) Pay? (\textbf{stīpendium, stīpendiī}.) A hired worker? (\textbf{mercēnārius, mercēnāriī}.) Toilsome? (\textbf{labōriōsus, labōriōsa, labōriōsum}.) Care? (\textbf{dīligentia, dīligentiae}.) To hire? (\textbf{condūcō, condūcere}.) To place? (\textbf{locō, locāre}.) A household servant? (\textbf{famulus, famulī}.) To do with energy? (\textbf{nāvō, nāvāre}.) To tire out? (\textbf{fatīgō, fatīgāre}.)
\end{tcolorbox}
